% from ega3-1-fr.tex
\oldpage[III]{98}
effet, lorsqu'il en est ainsi, tout $\mathcal O_X$-Module quasi-cohérent est
de la forme $\widetilde M$ (II, 2.7.7).
\end{env}

\begin{env}[2.1.7]
\label{III.2.1.7-fr}
Nous allons appliquer (\hyperref[III.2.1.5-fr]{2.1.5}) dans le cas
$S=A[T_0,\ldots,T_r]$, où $A$ est un anneau quelconque, les $T_i$ sont des
indéterminées, avec $M=S$ et $f_i=T_i$. On est donc essentiellement ramené
à calculer $H^\bullet((\mathbf T),S)$, où
$\mathbf T=(T_i)_{0\leq i\leq r}$.
\end{env}

\begin{lemma}[2.1.8]
\label{III.2.1.8-fr}
Si $S=A[T_0,\ldots,T_r]$, on a, avec
$\mathbf T=(T_i)_{0\leq i\leq r}$,
\begin{align}
\label{III.2.1.8.1-fr}
\tag{2.1.8.1}
H^i(\mathbf T^n,S)&=0 &&\text{si }i\neq r+1,\\
\label{III.2.1.8.2-fr}
\tag{2.1.8.2}
H^{r+1}(\mathbf T^n,S)&=S/(\mathbf T^n).
\end{align}
Le $A$-module $H^{r+1}(\mathbf T^n,S)$ a donc une base sur $A$ formée des
classes mod. $(\mathbf T^n)$ des monômes
\[
T_{\mathbf p}=T_0^{p_0}\cdots T_r^{p_r},
\qquad \mathbf p=(p_0,\ldots,p_r),\quad 0\leq p_i<n
\quad\text{pour tout }i.
\]
\end{lemma}

\begin{proof}
C'est une conséquence immédiate de (1.1.3.5) et de la prop. (1.1.4), dont
les hypothèses sont trivialement vérifiées.
\end{proof}

\begin{env}[2.1.9]
\label{III.2.1.9-fr}
Passons à la limite inductive sur $n$~; les relations
(\hyperref[III.2.1.8.1-fr]{2.1.8.1}) donnent
$H^i((\mathbf T),S)=0$ pour $i\neq r+1$. Pour $i=r+1$, le système inductif
est formé des $S/(\mathbf T^n)$, l'homomorphisme
\[
\varphi_{nm}:S/(\mathbf T^n)\longrightarrow S/(\mathbf T^m)
\qquad(0\leq n\leq m)
\]
étant la multiplication par $(T_0\cdots T_r)^{m-n}$. Pour
$n\geq\sup(p_i)_{0\leq i\leq r}$, désignons par
$\xi_{\mathbf p}^{(n)}=\xi_{p_0,\ldots,p_r}^{(n)}$ la classe de
$T_0^{n-p_0}\cdots T_r^{n-p_r}\pmod{(\mathbf T^n)}$~; on a alors
$\varphi_{nm}(\xi_{\mathbf p}^{(n)})=\xi_{\mathbf p}^{(m)}$, et ces éléments
ont donc même image canonique
$\xi_{\mathbf p}=\xi_{p_0,\ldots,p_r}$ dans la limite inductive
$H^{r+1}((\mathbf T),S)$~; en vertu de la définition du degré donnée dans
(\hyperref[III.2.1.2-fr]{2.1.2}), le degré de $\xi_{\mathbf p}$ est donc
égal à $-|\mathbf p|=-(p_0+p_1+\cdots+p_r)$. Il est clair que les
$\xi_{\mathbf p}^{(n)}$ pour $0<p_i\leq n$ et $0\leq i\leq r$ forment une
base de $S/(\mathbf T^n)$. On déduit donc aussitôt de
(\hyperref[III.2.1.8-fr]{2.1.8})~:
\end{env}

\begin{corollary}[2.1.10]
\label{III.2.1.10-fr}
Avec les notations de (\hyperref[III.2.1.8-fr]{2.1.8}), on a
\begin{equation}
\label{III.2.1.10.1-fr}
\tag{2.1.10.1}
H^i((\mathbf T),S)=0\qquad\text{pour }i\neq r+1,
\end{equation}
et $H^{r+1}((\mathbf T),S)$ est un $A$-module libre dont une base est formée
des éléments $\xi_{p_0,\ldots,p_r}$ tels que $p_i>0$ pour $0\leq i\leq r$.
\end{corollary}

\begin{env}[2.1.11]
\label{III.2.1.11-fr}
\emph{Remarque.} --- Soit $N$ un $A$-module quelconque et soit
$M=S\otimes_A N$~; le raisonnement de
(\hyperref[III.2.1.8-fr]{2.1.8}) montre que l'on a plus généralement
\begin{align}
\label{III.2.1.11.1-fr}
\tag{2.1.11.1}
H^i(\mathbf T^n,M)&=0 &&\text{si }i\neq r+1,\\
\label{III.2.1.11.2-fr}
\tag{2.1.11.2}
H^{r+1}(\mathbf T^n,M)&=(S/(\mathbf T^n))\otimes_A N,
\end{align}
car la dernière formule se déduit directement de (1.1.3.5), et d'autre
part il est clair que
\[
M/(T_0^nM+\cdots+T_{i-1}^nM)
=\bigl(S/(T_0^nS+\cdots+T_{i-1}^nS)\bigr)\otimes_A N,
\]
l'idéal $T_0^nS+\cdots+T_{i-1}^nS$ étant facteur direct dans le
$A$-module $S$~; cela permet d'appliquer (1.1.4) à $M$, et on obtient ainsi
(\hyperref[III.2.1.11.1-fr]{2.1.11.1}).
\end{env}

Combinant (\hyperref[III.2.1.10-fr]{2.1.10}) et
(\hyperref[III.2.1.5-fr]{2.1.5}), on obtient~:

\begin{proposition}[2.1.12]
\label{III.2.1.12-fr}
Soient $A$ un anneau, $r$ un entier $>0$, et
$X=\mathbf P_A^r$ (II, 4.1.1). Alors~:
\begin{enumerate}
\item[{\rm(i)}] On a $H^i(X,\mathcal O_X(*))=0$ pour $i\neq0,r$.
\item[{\rm(ii)}] L'homomorphisme canonique
$\alpha:S\to H^0(X,\mathcal O_X(*))$ (II, 2.6.2) est bijectif.
